\section*{Estatísticas}
\thispagestyle{empty} % Garante página sem rodapé
\addcontentsline{toc}{section}{Estatísticas}

% chama o Lua para gerar CSV
\directlua{
gerar_estatisticas_csv(
"../data/DadosBrutos_tmp.csv",
"../data/estatisticas_tmp.csv",
"../data/totais_tmp.csv",
"../data/auxiliares_tmp.csv",
"../data/estatisticas_mensal_tmp.csv"
)
}

%%% Carregar CSV
\DTLloaddb{mensal}{../data/estatisticas_mensal_tmp.csv}
\DTLloaddb{estatisticas}{../data/estatisticas_tmp.csv}
\DTLloaddb{auxiliares}{../data/auxiliares_tmp.csv}

%%% Escreve o texto inicial com as estatísticas gerais do período
\noindent\directlua{
texto_estatisticas_totais("../data/totais_tmp.csv")
tex.sprint(" ")
tipo_atividade("\detokenize{../data/DadosBrutos_tmp.csv}")
}

\noindent A tabela a seguir apresenta um resumo das atividades oficiais do clube durante o período
coberto por este documento. Foram considerados três indicadores: (i) a quantidade de \textbf{pranchetas} abertas; (ii) número de \textbf{participantes} conduzidos; e (iii) o total de \textbf{horas} em atividades oficiais na montanha.

\medskip
\begin{multicols}{2}

\begin{center}
\footnotesize
\begin{tabular}{rccr}
\textbf{Mês} & \textbf{Pranchetas} & \textbf{Participantes} & \textbf{Horas} \\
\hline
\DTLforeach{mensal}{
  \mes=Mes,
  \atividades=Atividades,
  \participantes=Participantes,
  \horas=Horas}{%
  \mes & \atividades & \participantes & \horas \\
}
\end{tabular}
\end{center}

\columnbreak

\begin{center}
\begin{tikzpicture}[scale=0.6]

  \begin{axis}[
    ybar, bar width=18pt,
    axis lines=left,
    ymin=0,
    symbolic x coords={2026-01,2026-02,2026-03},
    title={Pranchetas por Mês},
    xtick=data,
    enlarge x limits=0.1, enlarge y limits={upper, value=0.1}
    ]
    \addplot[
    fill=azulescuro,
    draw=azulescuro
    ] table[
    col sep=comma,
    x=Mes, y=Atividades
    ] {../data/estatisticas_mensal_tmp.csv};
  \end{axis}

\end{tikzpicture}
\end{center}

\end{multicols}

\subsection*{Guias e Auxiliares}
\noindent A seguir seguem as estatísticas dos guias do clube em pranchetas oficiais. A \textbf{Diretoria Técnica do CERJ} agradece e parabeniza a todos por dedicarem seu tempo para conduzir com segurança os sócios do CERJ.

\medskip

\begin{center}
\footnotesize
\begin{tabular}{rccr}
& \textbf{Pranchetas} & \textbf{Participantes} & \textbf{Horas} \\
\hline
\DTLforeach{estatisticas}{
  \g=Guia,
  \a=Atividades,
  \p=Participantes,
  \h=Horas}{%
\g & \a & \p & \h \\
}
\end{tabular}
\end{center}

\medskip

\noindent Além dos guias, muitos sócios, guias ou não, também dedicaram seu tempo a participar das atividades como \textbf{auxiliares de pranchetas}, ajudando os guias a organizar e conduzir os membros do clube nas montanhas. A \textbf{Diretoria Técnica do CERJ} agradece a todos e todas pela colaboração.

\medskip

\begin{center}
\footnotesize
\begin{tabular}{rc}
& \textbf{Pranchetas} \\
\hline
\DTLforeach{auxiliares}{
  \g=Auxiliar,
  \a=Atividades}{%
\g & \a \\
}
\end{tabular}
\end{center}
